\documentclass[leqno]{article}
% Shared mark scheme preamble for resources.danieljsmith.org: the pack
% preamble, unchanged, then the scheme's own layer. A scheme's .tex is
% rendered from its YAML (scripts/README.md, Mark schemes), never edited.
\input{\CurrentFilePath/pack_preamble.tex}

% The page: equal side margins, so the table runs the full width of the
% header rule. The footers' right edge follows the text's.
\geometry{left=0.8in, right=0.8in}
\makeatletter
\renewcommand{\djs@footR}[1]{%
  \put(\dimexpr1in+\hoffset+\oddsidemargin+\textwidth\relax,
    -\dimexpr\paperheight-0.32in\relax){\makebox[0pt][r]{#1}}}
\makeatother

\usepackage[export]{adjustbox}
\usepackage{tabularray}
\definecolor{msGrid}{RGB}{176, 172, 164}
\definecolor{msTint}{RGB}{245, 243, 239}

% The header on every page, the front page included, and the booklet label
% under the contents-page title. \djsQuestionPageStyle stays empty.
\newcommand{\djsSchemeHeader}[1]{%
  \djsPackHeader[markscheme]{\textit{Mark scheme} - #1}%
  \renewcommand{\djsPackTitle}{#1}%
  \renewcommand{\djsBookletLabel}{Mark scheme}%
}

% The front page: the contents heading, spaced as tocloft spaces the packs'
% "Questions", then the contents table. A question's name and page number
% both link to its first page.
\newcommand{\msContentsHeading}{%
  \vskip\cftbeforetoctitleskip
  {\noindent\cfttoctitlefont\contentsname\par}%
  \vskip\cftaftertoctitleskip}
\NewTblrEnviron{mscontents}
\SetTblrInner[mscontents]{
  colspec = {Q[wd=26mm, l] X[l] Q[wd=16mm, c] Q[wd=14mm, r]},
  width = \linewidth,
  rowsep = 4pt, colsep = 6pt,
  hlines = {0.4pt, msGrid}, vlines = {0.4pt, msGrid},
  hline{2} = {2.2pt, scarlet},
  row{1} = {font=\small\bfseries},
  row{Z} = {font=\bfseries},
  column{2} = {font=\small},
}
\newcommand{\msContentsQuestion}[1]{\hyperref[q:#1]{Question #1}}
\newcommand{\msContentsPage}[1]{\hyperref[q:#1]{\pageref*{q:#1}}}
\newcommand{\msSep}{\hspace{0.9em}}

% The conventions pages' layout: headings in the documents' own bold, tables
% ruled and padded as the scheme's, roomy lists.
\newcommand{\msHead}[1]{%
  \par\vspace{18pt}%
  {\noindent\bfseries #1\par}%
  \vspace{6pt}}
\NewTblrEnviron{mstblr}
\SetTblrInner[mstblr]{rowsep = 6pt, colsep = 6pt,
  hlines = {0.4pt, msGrid}, vlines = {0.4pt, msGrid}}
\newenvironment{msList}
  {\begin{itemize}[leftmargin=1.2em, itemsep=5pt, topsep=3pt, parsep=0pt]}
  {\end{itemize}}

% A question's first page: the link target and the PDF bookmark, as
% \questionitem gives the packs.
\newcommand{\msQnum}{}
\newcommand{\msQuestion}[1]{%
  \clearpage
  \renewcommand{\msQnum}{#1}%
  \phantomsection\label{q:#1}%
  \addcontentsline{toc}{subsection}{\protect\djsTocEntry{#1}}}

% The scheme's table: Question 16mm, Scheme the rest, Marks 16mm, Notes
% 44mm; every cell ruled in msGrid; the scarlet rule under the head row. A
% long table has no caption or "continued" foot; a continuation page opens
% with a muted line above the repeated head row. The rules across a part's
% rows, the alternative's tint and the total row are the rendered table's.
\NewTblrTheme{ms}{%
  \DefTblrTemplate{firsthead}{default}{}%
  \DefTblrTemplate{middlehead,lasthead}{default}{%
    {\small\itshape\color{djsMuted}Question \msQnum\ continued\par}%
    \vspace{6pt}}%
  \DefTblrTemplate{firstfoot,middlefoot,lastfoot}{default}{}%
}
\NewTblrEnviron{msscheme}
\SetTblrOuter[msscheme]{long, theme=ms, entry=none, label=none}
\SetTblrInner[msscheme]{
  colspec = {Q[wd=16mm, l] X[l] Q[wd=16mm, c] Q[wd=44mm, l]},
  width = \linewidth,
  rowsep = 6pt, colsep = 6pt,
  rowhead = 1,
  vlines = {0.4pt, msGrid},
  hline{1} = {0.4pt, msGrid}, hline{2} = {2.2pt, scarlet},
  hline{Z} = {0.4pt, msGrid},
  column{1} = {font=\bfseries}, column{3} = {font=\bfseries},
  column{4} = {font=\small},
  row{1} = {rowsep=4pt, font=\small\bfseries},
}

% Two marks on a row: each note starts with its code.
\newcommand{\msNote}[2]{\textbf{#1}\enspace #2}
% A full-row figure: the worked solution's own picture, scaled uniformly and
% only down, to the table less the Question column. The picture is drawn
% with the line width it has in the solutions document, at its list depth
% there (the optional argument, 1 by default: a part's working), so one
% sized by \linewidth keeps its size; tabularray sets the cell's \linewidth
% to the cell's measured width, which would shrink it at every pass. The
% widths are measured once, in a solution box at the packs' text width
% (pack_preamble.tex's margins, 0.8in and 1in).
\newlength{\msSpanW}
\makeatletter
\newcommand{\ms@keepLW}[1]{%
  \expandafter\xdef\csname ms@solnLW@#1\endcsname{\the\linewidth}}
\AtBeginDocument{%
  \setlength{\msSpanW}{\dimexpr\textwidth-16mm-28pt\relax}%
  \setbox\z@\vbox{%
    \begin{minipage}{\dimexpr\paperwidth-1.8in\relax}
      \begin{djsSolution}
        \ms@keepLW{0}%
        \begin{enumerate}\item \ms@keepLW{1}%
        \begin{enumerate}\item \ms@keepLW{2}%
        \begin{enumerate}\item \ms@keepLW{3}%
        \end{enumerate}\end{enumerate}\end{enumerate}
      \end{djsSolution}
    \end{minipage}}}
\newcommand{\msFigure}[2][1]{%
  \adjustbox{max width=\msSpanW, max totalheight=0.6\textheight, valign=t}{%
    \hsize=\csname ms@solnLW@#1\endcsname
    \linewidth=\hsize
    #2}}
\makeatother
% A figure mark's feature, in the Scheme column.
\newcommand{\msFeature}[1]{\textit{#1}}
% An alternative's title row.
\newcommand{\msAlternative}[1]{\textit{Alternative:} #1}
% Drafts only: a queried mark's code, the query's why under its note, and a
% flag in its part's place.
\newcommand{\msQueried}[2]{#1~\(\to\) #2?}
\newcommand{\msQuery}[1]{{\itshape\color{djsMuted}Query: #1}}
\newcommand{\msFlag}[2]{\textbf{#1:}\enspace #2}


\djsSchemeHeader{Systems of Differential Equations}
\djsSitePath{Core Pure}{Differential Equations}

\djsLastUpdated{25 October 2026}
\begin{document}
\djsSetQuestionCount{12}
\djsAddMarks{1}{10}
\djsAddMarks{2}{11}
\djsAddMarks{3}{11}
\djsAddMarks{4}{13}
\djsAddMarks{5}{10}
\djsAddMarks{6}{11}
\djsAddMarks{7}{18}
\djsAddMarks{8}{17}
\djsAddMarks{9}{16}
\djsAddMarks{10}{16}
\djsAddMarks{11}{15}
\djsAddMarks{12}{23}
\vspace*{6pt}
\djsContentsTitle
\vspace{30pt}
\msContentsHeading
\begin{mscontents}{}
Question & Marks by part & Marks & Page \\
\msContentsQuestion{1} & (a)\,3\msSep (b)\,7 & 10 & \msContentsPage{1} \\
\msContentsQuestion{2} & & 11 & \msContentsPage{2} \\
\msContentsQuestion{3} & & 11 & \msContentsPage{3} \\
\msContentsQuestion{4} & (a)\,2\msSep (b)\,11 & 13 & \msContentsPage{4} \\
\msContentsQuestion{5} & & 10 & \msContentsPage{5} \\
\msContentsQuestion{6} & & 11 & \msContentsPage{6} \\
\msContentsQuestion{7} & (a)\,3\msSep (b)\,6\msSep (c)\,2\msSep (d)\,6\msSep (e)\,1 & 18 & \msContentsPage{7} \\
\msContentsQuestion{8} & (a)\,3\msSep (b)\,4\msSep (c)\,3\msSep (d)\,7 & 17 & \msContentsPage{8} \\
\msContentsQuestion{9} & (a)\,2\msSep (b)\,3\msSep (c)\,11 & 16 & \msContentsPage{9} \\
\msContentsQuestion{10} & (a)\,1\msSep (b)\,3\msSep (c)(i)\,2\msSep (c)(ii)\,1\msSep (c)(iii)\,3\msSep (d)\,5\msSep (e)\,1 & 16 & \msContentsPage{10} \\
\msContentsQuestion{11} & (a)\,3\msSep (b)\,7\msSep (c)\,2\msSep (d)\,2\msSep (e)\,1 & 15 & \msContentsPage{11} \\
\msContentsQuestion{12} & (a)(i)\,5\msSep (a)(ii)\,2\msSep (b)(i)\,3\msSep (b)(ii)\,1\msSep (c)(i)\,2\msSep (c)(ii)\,4\msSep (c)(iii)\,3\msSep (d)(i)\,1\msSep (d)(ii)\,2 & 23 & \msContentsPage{12} \\
\SetCell[c=2]{r} Total & & 171 & \\
\end{mscontents}
\newpage
% The conventions pages, shared by every pack's mark scheme: one source of
% truth, \input after the front page, with no strand or pack logic. The keys
% on the first side, the rules on the second. The scheme layer defines
% \msHead, msList and mstblr.
\vspace*{6pt}%
{\noindent\cfttoctitlefont Reading the mark scheme\par}%
\vspace{10pt}%
These rules apply throughout, unless a note says otherwise.

\msHead{Marks}
Every mark is worth one, so \textbf{M1 A1} on a row is two marks. Only marks
coded ft follow through.\par
\vspace{10pt}%
\begin{mstblr}{
  colspec = {Q[wd=15mm, l] X[l]},
  width = \linewidth,
  hline{2} = {2.2pt, scarlet},
  column{1} = {font=\bfseries},
  row{1} = {font=\small\bfseries, rowsep=4pt},
}
Code & Meaning \\
M1 & Method: a correct method, applied to this question. A slip doesn't cost it if the method is clear. A formula must be used, not just quoted. \\
dM1 & Dependent method: as M1, but only if the method mark before it was earned, unless its note names another. \\
A1 & Accuracy: a correct answer or step. It needs the method mark it follows. \\
A1ft & As A1, following through from an earlier answer that the note names. \\
A1~AG & The last mark of a part whose answer is printed in the question: complete working, with no errors. \\
B1 & A correct value, fact or statement. It needs no method mark. \\
B1ft & As B1, following through from an earlier answer that the note names. \\
E1 & A written explanation, interpretation or conclusion. Concluding a proof by induction needs every earlier mark. \\
\end{mstblr}

\msHead{Abbreviations}
ft and AG qualify a code in the Marks column. The others appear in the
notes.\par
\vspace{10pt}%
\begin{mstblr}{
  colspec = {Q[wd=17mm, l] X[l] Q[wd=15mm, l] X[l]},
  width = \linewidth,
  column{1} = {font=\bfseries}, column{3} = {font=\bfseries},
}
ft & follow through & AG & answer given in the question \\
cao & correct answer only, with no follow-through & cso & correct solution only: no errors in the part \\
oe & or an equivalent form & awrt & any answer that rounds to this value \\
soi & seen, or implied by later working & isw & ignore working after a correct answer \\
SC & special case: credit outside the main method & their & the student's own earlier value \\
condone & a slip that doesn't cost the mark & [\;\,] & working the mark doesn't need \\
s.f. & significant figures & d.p. & decimal places \\
\end{mstblr}

\newpage
\msHead{Reading the scheme}
\begin{msList}
\item Each row is a line of working, with the marks it earns beside it.
  Where a row carries two marks, each note starts with its code.
\item \textbf{[\textit{n}]} ends each part with its total.
\item A shaded \textit{Alternative} is a second method for the same total.
  Mark an answer by the method it follows.
\item Any other valid method earns the equivalent marks, unless the question
  names a method or says \textit{hence}.
\item A sketch earns one B1 for each feature named below its figure.
\end{msList}

\msHead{Accuracy and answers}
\begin{msList}
\item Unless a question says otherwise, give answers to 3 significant
  figures, and angles in degrees to 1 decimal place. Angles in radians take
  3 significant figures.
\item A final answer with more figures is accepted if it rounds to the
  scheme's value; one with fewer is not. An answer spoiled by rounding too
  early loses its accuracy mark.
\item Where a question asks for an exact answer, a decimal loses the final
  accuracy mark. Otherwise either is accepted.
\item Wrong or missing units cost nothing unless the question asks for them.
\item A correct answer with no working earns full marks, unless the question
  needs the working: show that, prove, hence, a method it names, or an exact
  answer. A wrong answer with no working earns nothing. Calculator answers
  are treated the same way.
\item A correct answer, once reached, isn't spoiled by later working (isw).
\end{msList}

\msHead{Value of $\boldsymbol{g}$}
\begin{msList}
\item Take $g = 9.8\,\mathrm{m\,s^{-2}}$ unless a question states a value,
  which is then the only one accepted. Answers that use $g = 9.8$ may be
  given to 2 or 3 significant figures.
\item Where a question doesn't state $g$, an answer worked with $g = 10$
  throughout earns the same marks, keeping any value printed in the
  question. The worked solutions list the final answers with $g = 10$.
\end{msList}


\msQuestion{1}
\begin{msscheme}{
  hline{3-5,7-13} = {2-Z}{0.4pt, msGrid},
  hline{6,14} = {0.4pt, msGrid},
  row{14} = {font=\bfseries},
}
Question & Scheme & Marks & Notes \\
1(a) & {\(\frac{\mathrm{d}^2x}{\mathrm{d}t^2}+\frac{\mathrm{d}y}{\mathrm{d}t}=5\frac{\mathrm{d}x}{\mathrm{d}t}\)\\[4pt] \(\Rightarrow \frac{\mathrm{d}^2x}{\mathrm{d}t^2}+6x-2y=5\frac{\mathrm{d}x}{\mathrm{d}t}\)} & M1 & {Differentiates the first equation and replaces \(\frac{\mathrm{d}y}{\mathrm{d}t}\) using the second. Condone sign slips.} \\
& {\(\frac{\mathrm{d}^2x}{\mathrm{d}t^2}+6x-2\left(5x-\frac{\mathrm{d}x}{\mathrm{d}t}\right)=5\frac{\mathrm{d}x}{\mathrm{d}t}\)} & dM1 & {Substitutes \(y\) from the first equation to reach an equation in \(x\) and its derivatives only. Condone sign slips.} \\
& {\(\frac{\mathrm{d}^2x}{\mathrm{d}t^2}-3\frac{\mathrm{d}x}{\mathrm{d}t}-4x=0\)} & A1~AG & {Correct substituted equation leading without error to the printed result; collected steps accepted.} \\
\nopagebreak & & [3] & \\
1(b) & {\(\lambda^2-3\lambda-4=0 \Rightarrow \lambda=\ldots\)} & M1 & {Forms the auxiliary equation and solves it by any method, calculator allowed.} \\
& {\(\lambda=4,\ -1\)\\[4pt] \(x=Ae^{4t}+Be^{-t}\)} & A1 & \\
& {\(\frac{\mathrm{d}x}{\mathrm{d}t}=4Ae^{4t}-Be^{-t}\)\\[4pt] \(y=5\left(Ae^{4t}+Be^{-t}\right)-\left(4Ae^{4t}-Be^{-t}\right)\)} & M1 & {Differentiates their \(x\), of the form \(Ae^{pt}+Be^{qt}\), and substitutes into \(y=5x-\frac{\mathrm{d}x}{\mathrm{d}t}\). Condone sign slips.} \\
& {\(y=Ae^{4t}+6Be^{-t}\)} & A1 & \\
& {\(A+B=1,\quad A+6B=3\)} & M1 & {Substitutes \(t=0\), \(x=1\) and \(y=3\) into their \(x\) and \(y\) to form two equations in \(A\) and \(B\).} \\
& {\(\Rightarrow A=\ldots,\ B=\ldots\)} & dM1 & {Solves their two equations by any means, calculator allowed.} \\
& {\(x=\frac35e^{4t}+\frac25e^{-t}\)\\[4pt] \(y=\frac35e^{4t}+\frac{12}5e^{-t}\)} & A1 & {Both, oe.} \\
\nopagebreak & & [7] & \\
\nopagebreak \SetCell[c=2]{r} Total for Question 1 & & 10 & \\
\end{msscheme}

\msQuestion{2}
\begin{msscheme}{
  hline{3-12} = {2-Z}{0.4pt, msGrid},
  hline{13} = {0.4pt, msGrid},
  row{13} = {font=\bfseries},
}
Question & Scheme & Marks & Notes \\
2 & {\(y = -\frac{\mathrm{d}x}{\mathrm{d}t} \Rightarrow \frac{\mathrm{d}y}{\mathrm{d}t} = -\frac{\mathrm{d}^2x}{\mathrm{d}t^2}\)\\[4pt] \(\Rightarrow -\frac{\mathrm{d}^2x}{\mathrm{d}t^2} - x = 2e^t\)} & M1 & {Differentiates one equation and substitutes from the other to reach a second-order equation in \(x\) alone, or \(y\) alone. Condone sign slips.} \\
& {\(\frac{\mathrm{d}^2x}{\mathrm{d}t^2} + x = -2e^t\)} & A1 & {oe, or \(\frac{\mathrm{d}^2y}{\mathrm{d}t^2} + y = 2e^t\).} \\
& {\(t = 0: \frac{\mathrm{d}x}{\mathrm{d}t} + 1 = 0 \Rightarrow \frac{\mathrm{d}x}{\mathrm{d}t} = -1\)} & B1 & {Or \(\frac{\mathrm{d}y}{\mathrm{d}t} = 2\) at \(t = 0\) where \(x\) is eliminated.} \\
& {\(\lambda^2 + 1 = 0 \Rightarrow \lambda = \ldots\)} & M1 & {Forms the auxiliary equation of their second-order equation and solves it by any method.} \\
& {\(\lambda = \pm i\)\\[4pt] \(x = A\cos t + B\sin t\)} & A1 & {oe, in the variable of their equation.} \\
& {\(x = ae^t \Rightarrow ae^t + ae^t = -2e^t \Rightarrow a = \ldots\)} & M1 & {Substitutes a particular integral \(ae^t\) into their second-order equation and finds \(a\). Condone sign slips.} \\
& {\(x = A\cos t + B\sin t - e^t\)} & A1 & \\
& {\(A - 1 = 0\)\\[4pt] \(\frac{\mathrm{d}x}{\mathrm{d}t} = -A\sin t + B\cos t - e^t\)\\[4pt] \(\Rightarrow B - 1 = -1 \Rightarrow A = \ldots,\ B = \ldots\)} & M1 & {Uses both initial conditions in their general solution of the form \(A\cos t + B\sin t +\) their particular integral, and its attempted derivative, to find \(A\) and \(B\).} \\
& {\(x = \cos t - e^t\)} & A1 & \\
& {\(\frac{\mathrm{d}x}{\mathrm{d}t} = -\sin t - e^t \Rightarrow y = -\frac{\mathrm{d}x}{\mathrm{d}t} = \ldots\)} & M1 & {Differentiates their \(x\) and uses \(y = -\frac{\mathrm{d}x}{\mathrm{d}t}\), or any valid route to \(y\) leaving no undetermined constants.} \\
& {\(y = \sin t + e^t\)} & A1 & \\
\nopagebreak \SetCell[c=2]{r} Total for Question 2 & & 11 & \\
\end{msscheme}

\msQuestion{3}
\begin{msscheme}{
  hline{3-12} = {2-Z}{0.4pt, msGrid},
  hline{13} = {0.4pt, msGrid},
  row{13} = {font=\bfseries},
}
Question & Scheme & Marks & Notes \\
3 & {\(x = -\frac{\mathrm{d}y}{\mathrm{d}t} - y - 1 \Rightarrow \frac{\mathrm{d}x}{\mathrm{d}t} = -\frac{\mathrm{d}^2y}{\mathrm{d}t^2} - \frac{\mathrm{d}y}{\mathrm{d}t}\)\\[4pt] \(-\frac{\mathrm{d}^2y}{\mathrm{d}t^2} - \frac{\mathrm{d}y}{\mathrm{d}t}\)\\[4pt] \(\phantom{-} = 3\left(-\frac{\mathrm{d}y}{\mathrm{d}t} - y - 1\right) + 4y + 7\)} & M1 & {Makes \(x\) the subject of the second equation, differentiates and substitutes both into the first. Condone sign slips. Eliminating \(y\) instead, giving \(\frac{\mathrm{d}^2x}{\mathrm{d}t^2} - 2\frac{\mathrm{d}x}{\mathrm{d}t} + x = 3\), earns the corresponding marks.} \\
& {\(\frac{\mathrm{d}^2y}{\mathrm{d}t^2} - 2\frac{\mathrm{d}y}{\mathrm{d}t} + y = -4\)} & A1 & {oe} \\
& {\(\lambda^2 - 2\lambda + 1 = 0 \Rightarrow \lambda = \ldots\)} & M1 & {Forms the auxiliary equation of their second-order equation and solves it by any means.} \\
& {\(\lambda = 1\) (repeated)\\[4pt] \((y =)\ (A + Bt)e^t\)} & A1 & \\
& {\(y = k \Rightarrow k = -4\)} & B1 & \\
& {\(y = (A + Bt)e^t - 4\)} & A1ft & {ft their complementary function plus their constant particular integral; must be \(y = \ldots\).} \\
& {\(y(0) = -3 \Rightarrow A - 4 = -3 \Rightarrow A = \ldots\ [= 1]\)} & M1 & {Substitutes \(t = 0\), \(y = -3\) into their general solution to find \(A\).} \\
& {\(\frac{\mathrm{d}y}{\mathrm{d}t} = (A + Bt + B)e^t\)\\[4pt] \(x = -(A + Bt + B)e^t - (A + Bt)e^t + 4 - 1\)} & M1 & {Differentiates their \(y\), with the product rule on \((A + Bt)e^t\), and substitutes into \(x\) made the subject of the second equation, or an equivalent relation from the system.} \\
& {\(x = -(2A + 2Bt + B)e^t + 3\)} & A1 & {oe, in terms of \(A\) and \(B\) or with \(A = 1\).} \\
& {\(x(0) = 5 \Rightarrow -(2 + B) + 3 = 5 \Rightarrow B = \ldots\)} & M1 & {Substitutes \(t = 0\), \(x = 5\) into their \(x\), or uses \(\frac{\mathrm{d}y}{\mathrm{d}t} = -3\) at \(t = 0\), to find \(B\).} \\
& {\(x = (2 + 8t)e^t + 3\)\\[4pt] \(y = (1 - 4t)e^t - 4\)} & A1 & {Both, oe.} \\
\nopagebreak \SetCell[c=2]{r} Total for Question 3 & & 11 & \\
\end{msscheme}

\msQuestion{4}
\begin{msscheme}{
  hline{3-4,6-16} = {2-Z}{0.4pt, msGrid},
  hline{5,17} = {0.4pt, msGrid},
  row{17} = {font=\bfseries},
}
Question & Scheme & Marks & Notes \\
4(a) & {\(x+2y-7=0,\ 2x+y-8=0 \Rightarrow x = \ldots,\ y = \ldots\)} & M1 & {Sets both expressions equal to zero and solves the simultaneous equations by any method.} \\
& {\((3,\ 2)\)} & A1 & \\
\nopagebreak & & [2] & \\
4(b) & {\(\frac{\mathrm{d}^2x}{\mathrm{d}t^2} = \frac{\mathrm{d}x}{\mathrm{d}t} + 2\frac{\mathrm{d}y}{\mathrm{d}t}\)\\[4pt] \(\phantom{\frac{\mathrm{d}^2x}{\mathrm{d}t^2}} = \frac{\mathrm{d}x}{\mathrm{d}t} + 4x + 2y - 16\)\\[4pt] with \(2y = \frac{\mathrm{d}x}{\mathrm{d}t} - x + 7\)} & M1 & {Differentiates one equation, substitutes the other and eliminates one variable to reach a second-order equation in a single variable. Condone sign slips.} \\
& {\(\frac{\mathrm{d}^2x}{\mathrm{d}t^2} - 2\frac{\mathrm{d}x}{\mathrm{d}t} - 3x = -9\)} & A1 & {oe} \\
& {\(\lambda^2 - 2\lambda - 3 = 0 \Rightarrow \lambda = \ldots\)} & M1 & {Forms the auxiliary equation of their second-order equation and solves it by factorising, formula or completing the square; calculator allowed.} \\
& {\(Ae^{3t} + Be^{-t}\)} & A1 & \\
& {\(-3k = -9 \Rightarrow k = 3\)} & B1 & \\
& {\(x = Ae^{3t} + Be^{-t} + 3\)} & A1ft & {ft their complementary function plus their constant particular integral.} \\
& {\(A + B + 3 = 4\)\\[4pt] \(3A - B = 4 + 2(0) - 7\)} & M1 & {Forms two equations in \(A\) and \(B\): \(x(0) = 4\) and a second condition from the system at \(t = 0\), e.g. \(\frac{\mathrm{d}x}{\mathrm{d}t} = -3\) in their differentiated solution. Condone slips.} \\
& {\(\Rightarrow A = \ldots,\ B = \ldots\)} & dM1 & {Solves their two equations for \(A\) and \(B\); calculator allowed.} \\
& {\(x = 3 - \frac12 e^{3t} + \frac32 e^{-t}\)} & A1 & {oe} \\
& {\(2y = \their{-\frac32 e^{3t} - \frac32 e^{-t}}\)\\[4pt] \(\phantom{2y =} - \their{3 - \frac12 e^{3t} + \frac32 e^{-t}} + 7\)} & M1 & {Substitutes their \(x\) and their \(\frac{\mathrm{d}x}{\mathrm{d}t}\) into the first equation to find \(y\), with their \(x\) of the form \(pe^{\alpha t} + qe^{\beta t} + c\), numerical constants.} \\
& {\(y = 2 - \frac12 e^{3t} - \frac32 e^{-t}\)} & A1 & {oe} \\
\nopagebreak & & [11] & \\
\nopagebreak \SetCell[c=2]{r} Total for Question 4 & & 13 & \\
\end{msscheme}

\msQuestion{5}
\begin{msscheme}{
  hline{3-11} = {2-Z}{0.4pt, msGrid},
  hline{12} = {0.4pt, msGrid},
  row{12} = {font=\bfseries},
}
Question & Scheme & Marks & Notes \\
5 & {\(y = 2x - \frac{\mathrm{d}x}{\mathrm{d}t} \Rightarrow \frac{\mathrm{d}y}{\mathrm{d}t} = 2\frac{\mathrm{d}x}{\mathrm{d}t} - \frac{\mathrm{d}^2x}{\mathrm{d}t^2}\)\\[4pt] \(2\frac{\mathrm{d}x}{\mathrm{d}t} - \frac{\mathrm{d}^2x}{\mathrm{d}t^2} = x + 2\left(2x - \frac{\mathrm{d}x}{\mathrm{d}t}\right)\)} & M1 & {Makes one variable the subject of one equation, differentiates it and substitutes into the other, eliminating that variable. Condone sign slips.} \\
& {\(\frac{\mathrm{d}^2x}{\mathrm{d}t^2} - 4\frac{\mathrm{d}x}{\mathrm{d}t} + 5x = 0\)} & A1 & {oe, including the same equation in \(y\).} \\
& {\(\lambda^2 - 4\lambda + 5 = 0 \Rightarrow \lambda = \ldots\)} & M1 & {Forms the auxiliary equation of their second-order equation and solves it by any means, calculator allowed.} \\
& {\([\lambda = 2 \pm i]\)\\[4pt] \(x = e^{2t}(A\cos t + B\sin t)\)} & A1 & {Any letters for the constants; must be in terms of \(t\).} \\
& {\(t = 0: \frac{\mathrm{d}x}{\mathrm{d}t} = 2(1) - 2 = 0\)} & B1 & \\
& {\(\frac{\mathrm{d}x}{\mathrm{d}t} = 2e^{2t}(A\cos t + B\sin t)\)\\[4pt] \(\phantom{\frac{\mathrm{d}x}{\mathrm{d}t}} + e^{2t}(-A\sin t + B\cos t)\)} & M1 & {Differentiates their complementary function by the product rule, with both terms present. Condone sign slips in differentiating \(\sin t\) and \(\cos t\).} \\
& {\(A = 1,\ 2A + B = 0 \Rightarrow B = \ldots\)} & dM1 & {Substitutes \(t = 0\) into their \(x\) and their derivative, using \(x = 1\) and their value of \(\frac{\mathrm{d}x}{\mathrm{d}t}\), and solves for both constants.} \\
& {\(f(t) = e^{2t}(\cos t - 2\sin t)\)} & A1 & {oe simplified, such as \(e^{2t}\cos t - 2e^{2t}\sin t\).} \\
& {\(y = 2e^{2t}(\cos t - 2\sin t) + 5e^{2t}\sin t\)} & M1 & {Finds the second function from their first using one of the given equations, or by solving its own second-order equation with both its initial conditions.} \\
& {\(g(t) = e^{2t}(2\cos t + \sin t)\)} & A1 & {oe simplified, such as \(2e^{2t}\cos t + e^{2t}\sin t\).} \\
\nopagebreak \SetCell[c=2]{r} Total for Question 5 & & 10 & \\
\end{msscheme}

\msQuestion{6}
\begin{msscheme}{
  hline{3-12} = {2-Z}{0.4pt, msGrid},
  hline{13} = {0.4pt, msGrid},
  row{13} = {font=\bfseries},
}
Question & Scheme & Marks & Notes \\
6 & {\(y = -\frac{\mathrm{d}x}{\mathrm{d}t} \Rightarrow \frac{\mathrm{d}y}{\mathrm{d}t} = -\frac{\mathrm{d}^2x}{\mathrm{d}t^2}\)\\[4pt] \(\phantom{y} \Rightarrow -\frac{\mathrm{d}^2x}{\mathrm{d}t^2} - x = 2\)} & M1 & {Makes one variable the subject of one equation, differentiates and substitutes into the other, reaching a second-order equation in one variable. Condone sign slips.} \\
& {\(\frac{\mathrm{d}^2x}{\mathrm{d}t^2} + x = -2\)} & A1 & {oe} \\
& {\(\lambda^2 + 1 = 0 \Rightarrow \lambda = \ldots\)} & M1 & {Forms the auxiliary equation of their second-order equation and solves it by any means.} \\
& {\(\lambda = \pm i\)\\[4pt] \((x =)\ A\cos t + B\sin t\)} & A1 & \\
& {\(x = k \Rightarrow k = -2\)} & B1 & \\
& {\(x = A\cos t + B\sin t - 2\)} & A1ft & {ft their complementary function plus their constant particular integral, as \(x = \ldots\) with two arbitrary constants.} \\
& {\(t = \frac{\pi}{2}\): \(\frac{\mathrm{d}x}{\mathrm{d}t} = -y = -2\)} & B1 & {soi, e.g. by \(-A = -2\).} \\
& {\(1 = A\cos\frac{\pi}{2} + B\sin\frac{\pi}{2} - 2\)\\[4pt] \(\their{{-2}} = -A\sin\frac{\pi}{2} + B\cos\frac{\pi}{2}\)\\[4pt] \(\phantom{A} \Rightarrow A = \ldots,\ B = \ldots\)} & M1 & {Substitutes \(t = \frac{\pi}{2}\), \(x = 1\) into their general solution and their derivative value into its derivative, finding both constants. Condone differentiation sign slips.} \\
& {\(x = 2\cos t + 3\sin t - 2\)} & A1 & \\
& {\(y = -\frac{\mathrm{d}x}{\mathrm{d}t} = -(-2\sin t + 3\cos t)\)} & M1 & {Differentiates their \(x\) and uses \(y = -\frac{\mathrm{d}x}{\mathrm{d}t}\), or another relation from the system, without new arbitrary constants. Condone sign slips.} \\
& {\(y = 2\sin t - 3\cos t\)} & A1 & \\
\nopagebreak \SetCell[c=2]{r} Total for Question 6 & & 11 & \\
\end{msscheme}

\msQuestion{7}
\begin{msscheme}{
  hline{3-5,7-12,14-15,17-22,24} = {2-Z}{0.4pt, msGrid},
  hline{6,13,16,23,25} = {0.4pt, msGrid},
  row{25} = {font=\bfseries},
}
Question & Scheme & Marks & Notes \\
7(a) & {\(\frac{\mathrm{d}^2w}{\mathrm{d}t^2}=4\frac{\mathrm{d}w}{\mathrm{d}t}-\frac{\mathrm{d}s}{\mathrm{d}t}\)\\[4pt] \(\phantom{\frac{\mathrm{d}^2w}{\mathrm{d}t^2}}=4\frac{\mathrm{d}w}{\mathrm{d}t}-\left(6w-s-48e^{-t}\right)\)} & M1 & {Differentiates the first equation and uses the second to replace \(\frac{\mathrm{d}s}{\mathrm{d}t}\), or equivalent. Condone sign slips.} \\
& {\(\phantom{\frac{\mathrm{d}^2w}{\mathrm{d}t^2}}=4\frac{\mathrm{d}w}{\mathrm{d}t}-6w\)\\[4pt] \(\phantom{\frac{\mathrm{d}^2w}{\mathrm{d}t^2}=}+\left(4w-\frac{\mathrm{d}w}{\mathrm{d}t}\right)+48e^{-t}\)} & dM1 & {Eliminates \(s\) using \(s=4w-\frac{\mathrm{d}w}{\mathrm{d}t}\), reaching an equation in \(w\) and its derivatives only.} \\
& {\(\frac{\mathrm{d}^2w}{\mathrm{d}t^2}-3\frac{\mathrm{d}w}{\mathrm{d}t}+2w=48e^{-t}\)} & A1~AG & {Correct substituted equation leading without error to the printed result; collected steps accepted.} \\
\nopagebreak & & [3] & \\
7(b) & {\(\lambda^2-3\lambda+2=0 \Rightarrow \lambda=\ldots\)} & M1 & {Forms the auxiliary equation and solves it by factorising with product of constants 2, the formula or completing the square; calculator accepted.} \\
& {\(\lambda=1,\ 2\)\\[4pt] \((w=)\,Ae^{t}+Be^{2t}\)} & A1 & \\
& {\(w=ae^{-t}\)} & B1 & \\
& {\(ae^{-t}+3ae^{-t}+2ae^{-t}=48e^{-t} \Rightarrow a=\ldots\)} & M1 & {Differentiates their trial form twice, substitutes into the equation and equates coefficients to find \(a\). Condone sign slips.} \\
& {\(8e^{-t}\)} & A1 & \\
& {\(w=Ae^{t}+Be^{2t}+8e^{-t}\)} & A1ft & {ft their complementary function of the form \(Ae^{\alpha t}+Be^{\beta t}\) plus their particular integral. Must be \(w=\ldots\) in terms of \(t\).} \\
\nopagebreak & & [6] & \\
7(c) & {\(\frac{\mathrm{d}w}{\mathrm{d}t}=Ae^{t}+2Be^{2t}-8e^{-t}\)\\[4pt] \(s=4(\ldots)-(\ldots)\)} & M1 & {Differentiates their \(w\) and substitutes into \(s=4w-\frac{\mathrm{d}w}{\mathrm{d}t}\). Condone coefficient slips.} \\
& {\(s=3Ae^{t}+2Be^{2t}+40e^{-t}\)} & A1 & {Constants must match those in \(w\).} \\
\nopagebreak & & [2] & \\
7(d) & {\(30=A+B+8\), \(110=3A+2B+40\)} & M1 & {Substitutes \(t=0\), \(w=30\), \(s=110\) into their \(w\) and \(s\), forming two equations in \(A\) and \(B\).} \\
& {\(A+B=22,\ 3A+2B=70 \Rightarrow A=\ldots,\ B=\ldots\)} & dM1 & {Solves their two linear equations by any method; calculator accepted.} \\
& {\(A=26,\ B=-4\)} & A1 & \\
& {\(26e^{T}-4e^{2T}+8e^{-T}=0\)\\[4pt] \(\Rightarrow 2x^3-13x^2-4=0,\ x=e^{T}\)} & M1 & {Sets their \(w=0\) at \(t=T\) and forms a cubic in \(e^{T}\), or solves the equation in \(T\) directly. Need not be simplified.} \\
& {\(x=6.547\ldots\)} & B1 & {awrt 6.55; implied by a correct \(T\).} \\
& {\(T=\ln x=1.879\)} & A1 & \\
\nopagebreak & & [6] & \\
7(e) & {The model predicts negative values of \(w\) after \(T\)} & E1 & {States a limitation specific to this model, e.g. negative \(w\) or \(s\) predicted, \(s\) not an integer, unbounded growth. Not for ``not accurate'' alone.} \\
\nopagebreak & & [1] & \\
\nopagebreak \SetCell[c=2]{r} Total for Question 7 & & 18 & \\
\end{msscheme}

\msQuestion{8}
\begin{msscheme}{
  hline{3-5,7-9,11-13,15-21} = {2-Z}{0.4pt, msGrid},
  hline{6,10,14,22} = {0.4pt, msGrid},
  row{22} = {font=\bfseries},
}
Question & Scheme & Marks & Notes \\
8(a) & {\(\frac{\mathrm{d}^2s}{\mathrm{d}t^2} = -\frac{1}{20}\frac{\mathrm{d}m}{\mathrm{d}t} - \frac{1}{10}\frac{\mathrm{d}s}{\mathrm{d}t}\)} & M1 & {Differentiates the \(\frac{\mathrm{d}s}{\mathrm{d}t}\) equation, or \(m\) made its subject, with respect to \(t\). Condone coefficient slips.} \\
& {\(\phantom{\frac{\mathrm{d}^2s}{\mathrm{d}t^2}} = -\frac{1}{20}\left(\frac14 s\right) - \frac{1}{10}\frac{\mathrm{d}s}{\mathrm{d}t}\)} & dM1 & {Replaces \(\frac{\mathrm{d}m}{\mathrm{d}t}\) by \(\frac14 s\) to eliminate \(m\).} \\
& {\(\frac{\mathrm{d}^2s}{\mathrm{d}t^2} + \frac{1}{10}\frac{\mathrm{d}s}{\mathrm{d}t} + \frac{1}{80}s = 0\)} & A1~AG & {Correct substituted equation leading without error to the printed result; combined rearrangement accepted.} \\
\nopagebreak & & [3] & \\
8(b) & {\(\lambda^2 + \frac{1}{10}\lambda + \frac{1}{80} = 0 \Rightarrow \lambda = \ldots\)} & M1 & {Forms the auxiliary equation and solves it by formula with their values, completing the square or calculator.} \\
& {\(\lambda = -\frac{1}{20} \pm \frac{1}{10}i\)} & A1 & \\
& {\(s = e^{-t/20}\left(A\cos\frac{t}{10} + B\sin\frac{t}{10}\right)\)} & dM1 A1 & {\msNote{dM1}{Uses \(e^{pt}(A\cos qt + B\sin qt)\) for their roots \(p \pm qi\); a complex exponential form does not score.}} \\
\nopagebreak & & [4] & \\
8(c) & {\(m = -20\left(\frac{\mathrm{d}s}{\mathrm{d}t} + \frac{1}{10}s\right)\)\\[4pt] \(\frac{\mathrm{d}s}{\mathrm{d}t} = -\frac{1}{20}e^{-t/20}\left(A\cos\frac{t}{10} + B\sin\frac{t}{10}\right)\)\\[4pt] \(\phantom{\frac{\mathrm{d}s}{\mathrm{d}t} =} + e^{-t/20}\left(\frac{B}{10}\cos\frac{t}{10} - \frac{A}{10}\sin\frac{t}{10}\right)\)} & M1 & {Makes \(m\) the subject of the second equation and differentiates their \(s\), of the form \(e^{pt}(A\cos qt + B\sin qt)\), by the product rule. Condone sign and coefficient slips.} \\
& {\(m = -20e^{-t/20}\Big[\left(\frac{A}{20} + \frac{B}{10}\right)\cos\frac{t}{10}\)\\[4pt] \(\phantom{m = -20e^{-t/20}\Big[} + \left(\frac{B}{20} - \frac{A}{10}\right)\sin\frac{t}{10}\Big]\)} & dM1 & {Substitutes their \(s\) and \(\frac{\mathrm{d}s}{\mathrm{d}t}\) into their expression for \(m\). Need not be simplified.} \\
& {\(m = e^{-t/20}\Big[(-A-2B)\cos\frac{t}{10}\)\\[4pt] \(\phantom{m = e^{-t/20}\Big[} + (2A-B)\sin\frac{t}{10}\Big]\)} & A1 & {oe, using the same constants as \(s\).} \\
\nopagebreak & & [3] & \\
8(d) & {\(s(0) = A = 20\)\\[4pt] \(m(0) = -A - 2B = 20 \Rightarrow B = \ldots\)} & M1 & {Substitutes \(t = 0\), \(s = 20\), \(m = 20\) into their \(s\) and \(m\) and solves for both constants.} \\
& {\(A = 20,\ B = -20\)} & A1 & \\
& {\(20e^{-t/20}\left(\cos\frac{t}{10} - \sin\frac{t}{10}\right) = 0\)\\[4pt] \(\Rightarrow \tan\frac{t}{10} = 1 \Rightarrow t = \ldots\)} & M1 & {Sets their \(s = 0\) and solves for \(t\), e.g. via \(\tan\frac{t}{10} = k\) or a harmonic form.} \\
& {\(t = \frac{5\pi}{2} \approx 7.85\); 2028} & A1 & {\(\frac{5\pi}{2}\) or awrt 7.85, the first positive root, and 2028.} \\
& {\(m = 20e^{-\pi/8}\left(\cos\frac{\pi}{4} + 3\sin\frac{\pi}{4}\right)\)} & M1 & {Substitutes their \(A\), \(B\) and their \(t\) into their \(m\).} \\
& {\(m = 40\sqrt{2}\,e^{-\pi/8} \approx 38.2\)} & A1 & {awrt 38.2, or 38, or exact equivalent.} \\
& {\(s < 0\) for \(t > \frac{5\pi}{2}\): impossible\\[4pt] No saplings and 38 trees, fewer than 40\\[4pt] \(\Rightarrow\) unsuitable in the long term} & E1 & {Concludes the model is unsuitable, at least long term, citing their \(t\) or \(m\): e.g. \(s < 0\) afterwards, or no saplings with fewer than the original 40 trees.} \\
\nopagebreak & & [7] & \\
\nopagebreak \SetCell[c=2]{r} Total for Question 8 & & 17 & \\
\end{msscheme}

\msQuestion{9}
\begin{msscheme}{
  hline{3-4,6-8,10-20} = {2-Z}{0.4pt, msGrid},
  hline{5,9,21} = {0.4pt, msGrid},
  row{21} = {font=\bfseries},
}
Question & Scheme & Marks & Notes \\
9(a) & {\(\frac{\mathrm{d}x}{\mathrm{d}t} = \frac{r}{500}y - \frac{r}{300}x\)\\[4pt] \(\Rightarrow \frac{r}{300} = \frac{1}{20}\) or \(\frac{r}{500} = \frac{3}{100}\)} & M1 & {Forms the salt balance for \(A\) with inflow \(\frac{ry}{500}\) and an outflow term in \(\frac{x}{300}\), and equates a coefficient with the given equation.} \\
& {\(r = 15\)} & A1 & \\
\nopagebreak & & [2] & \\
9(b) & {Inflow to \(B\) \(= 15 - 3 = 12\)} & B1 & \\
& {\(\frac{\mathrm{d}y}{\mathrm{d}t} = 12\mu + \frac{3x}{300} - \frac{15y}{500}\)\\[4pt] \(\Rightarrow 12\mu = 24\)} & M1 & {Forms the salt balance for \(B\): \(\mu\) times their inflow rate, \(\frac{3x}{300}\) and an outflow term in \(\frac{y}{500}\); equates the constant term with 24.} \\
& {\(\mu = 2\)} & A1~AG & {Correct balance with rates 12 and 15, compared with the given equation, leading without error to \(\mu = 2\).} \\
\nopagebreak & & [3] & \\
9(c) & {\(\frac{\mathrm{d}^2x}{\mathrm{d}t^2} = -\frac{1}{20}\frac{\mathrm{d}x}{\mathrm{d}t}\)\\[4pt] \(\phantom{\frac{\mathrm{d}^2x}{\mathrm{d}t^2}} + \frac{3}{100}\left(\frac{x}{100} - \frac{3y}{100} + 24\right)\)\\[4pt] with \(3y = 100\frac{\mathrm{d}x}{\mathrm{d}t} + 5x\)} & M1 & {Differentiates one equation, substitutes the other and eliminates \(y\) (or \(x\)) to reach a second-order equation in one variable. Condone coefficient slips.} \\
& {\(\frac{\mathrm{d}^2x}{\mathrm{d}t^2} + \frac{2}{25}\frac{\mathrm{d}x}{\mathrm{d}t} + \frac{3}{2500}x = \frac{18}{25}\)} & A1 & {oe, such as \(2500\frac{\mathrm{d}^2x}{\mathrm{d}t^2} + 200\frac{\mathrm{d}x}{\mathrm{d}t} + 3x = 1800\).} \\
& {\(\lambda^2 + \frac{2}{25}\lambda + \frac{3}{2500} = 0 \Rightarrow \lambda = \ldots\)} & M1 & {Forms the auxiliary equation from their second-order equation and solves it by factorising, formula, completing the square or calculator.} \\
& {\(\lambda = -\frac{1}{50},\ -\frac{3}{50}\)\\[4pt] \((x =)\ Ae^{-t/50} + Be^{-3t/50}\)} & A1 & \\
& {\(\frac{3}{2500}k = \frac{18}{25} \Rightarrow k = 600\)} & B1 & \\
& {\(x = 600 + Ae^{-t/50} + Be^{-3t/50}\)} & A1ft & {ft their complementary function plus their constant particular integral, with two arbitrary constants.} \\
& {\(t = 0\): \(\frac{\mathrm{d}x}{\mathrm{d}t} = -\frac{1}{20}(300) + \frac{3}{100}(1200) = 21\)} & B1 & \\
& {\(600 + A + B = 300\)\\[4pt] \(-\frac{1}{50}A - \frac{3}{50}B = \their{21}\)\\[4pt] \(\Rightarrow A = \ldots,\ B = \ldots\)} & M1 & {Uses \(x = 300\) and their \(\frac{\mathrm{d}x}{\mathrm{d}t}\) (or \(y = 1200\)) at \(t = 0\) in their general solution and its derivative, and solves for both constants. Condone slips.} \\
& {\(x = 600 + 75e^{-t/50} - 375e^{-3t/50}\)} & A1 & \\
& {\(3y = 100\frac{\mathrm{d}x}{\mathrm{d}t} + 5x\)\\[4pt] \(\phantom{3y} = 100\left(-\tfrac32 e^{-t/50} + \tfrac{45}{2}e^{-3t/50}\right) + 5x\)} & M1 & {Substitutes their \(x\) and its derivative into the first equation rearranged for \(y\). Condone slips in differentiating exponentials.} \\
& {\(y = 1000 + 75e^{-t/50} + 125e^{-3t/50}\)} & A1 & {oe} \\
\nopagebreak & & [11] & \\
\nopagebreak \SetCell[c=2]{r} Total for Question 9 & & 16 & \\
\end{msscheme}

\msQuestion{10}
\begin{msscheme}{
  hline{3,5-7,9-10,12,14-16,18-22,24} = {2-Z}{0.4pt, msGrid},
  hline{4,8,11,13,17,23,25} = {0.4pt, msGrid},
  row{25} = {font=\bfseries},
}
Question & Scheme & Marks & Notes \\
10(a) & {If \(B > A\), water is pumped into tank A, so \(k(B-A) > 0\)\\[4pt] \(\Rightarrow k > 0\)} & E1 & {Links the direction of pumping to the sign of \(k(B-A)\) to conclude \(k > 0\). One case, \(B > A\) or \(A > B\), suffices.} \\
\nopagebreak & & [1] & \\
10(b) & {\(\frac{\mathrm{d}^2A}{\mathrm{d}t^2} = \frac{\mathrm{d}B}{\mathrm{d}t} - \frac{\mathrm{d}A}{\mathrm{d}t} + 8t\)\\[4pt] \(\phantom{\frac{\mathrm{d}^2A}{\mathrm{d}t^2}} = A - B - \frac{\mathrm{d}A}{\mathrm{d}t} + 8t\)} & M1 & {Differentiates the first equation with respect to \(t\) and substitutes for \(\frac{\mathrm{d}B}{\mathrm{d}t}\). Condone a slip in differentiating \(4t^2\).} \\
& {\(A - B = 4t^2 - \frac{\mathrm{d}A}{\mathrm{d}t}\)\\[4pt] \(\Rightarrow \frac{\mathrm{d}^2A}{\mathrm{d}t^2} = \left(4t^2 - \frac{\mathrm{d}A}{\mathrm{d}t}\right) - \frac{\mathrm{d}A}{\mathrm{d}t} + 8t\)} & dM1 & {Uses the first equation to eliminate \(B\), leaving only \(A\), its derivatives and \(t\). Condone sign slips.} \\
& {\(\frac{\mathrm{d}^2A}{\mathrm{d}t^2} + 2\frac{\mathrm{d}A}{\mathrm{d}t} = 4t^2 + 8t\)} & A1~AG & {Correct substituted equations leading without error to the printed equation; collected steps accepted.} \\
\nopagebreak & & [3] & \\
10(c)(i) & {\(\lambda^2 + 2\lambda = 0 \Rightarrow \lambda = \ldots\)} & M1 & {Forms the auxiliary equation and solves it by any means.} \\
& {\((A =)\ C_1 + C_2e^{-2t}\)} & A1 & {oe, any two arbitrary constants.} \\
\nopagebreak & & [2] & \\
10(c)(ii) & {\(\frac{\mathrm{d}^2A}{\mathrm{d}t^2} + 2\frac{\mathrm{d}A}{\mathrm{d}t} = 4at + 2a + 2b\): no \(t^2\) term\\[4pt] (the constant \(c\) is in the complementary function)} & E1 & {Shows or states that the trial gives no \(t^2\) term, or that its constant term is in the complementary function. Either reason suffices.} \\
\nopagebreak & & [1] & \\
10(c)(iii) & {\((6at + 2b) + 2(3at^2 + 2bt + c) = 4t^2 + 8t\)\\[4pt] \(\Rightarrow 6a = 4,\ 6a + 4b = 8,\ 2b + 2c = 0\)} & M1 & {Differentiates the trial twice, substitutes and equates coefficients of at least two powers of \(t\). Condone slips in differentiation.} \\
& {\(a = \frac23,\ b = 1,\ c = -1\)} & A1 & \\
& {\(A = C_1 + C_2e^{-2t} + \frac23t^3 + t^2 - t\)} & A1ft & {ft their complementary function and their \(a\), \(b\), \(c\); needs \(A =\) and two arbitrary constants.} \\
\nopagebreak & & [3] & \\
10(d) & {\(\frac{\mathrm{d}A}{\mathrm{d}t} = -2C_2e^{-2t} + 2t^2 + 2t - 1\)\\[4pt] \(t = 0:\ -2C_2 - 1 = 13\)} & M1 & {Differentiates their general solution, of the form \(C_1 + C_2e^{-2t} +\) polynomial, and substitutes \(t = 0\), \(\frac{\mathrm{d}A}{\mathrm{d}t} = 13\). Condone differentiation slips.} \\
& {\(t = 0:\ C_1 + C_2 = 52\)} & M1 & {Substitutes \(t = 0\), \(A = 52\) into their general solution.} \\
& {\(C_2 = -7,\ C_1 = 59\)} & A1 & \\
& {\(A = \their{59 - 7e^{-1}} + \frac23(0.5)^3 + 0.5^2 - 0.5\)} & M1 & {Substitutes \(t = 0.5\), not \(t = 30\), into their solution with their constants.} \\
& {\(A \approx 56.3\ (\mathrm{m}^3)\)} & A1 & {awrt 56.3, or \(59 - 7e^{-1} - \frac16\) oe.} \\
\nopagebreak & & [5] & \\
10(e) & {\(4t^2\), and so \(A\), grows without bound, but the tanks have finite capacity} & E1 & {Refers to the unbounded inflow or volume and a physical limit such as tank capacity. Not for 'the model is a simplification' alone.} \\
\nopagebreak & & [1] & \\
\nopagebreak \SetCell[c=2]{r} Total for Question 10 & & 16 & \\
\end{msscheme}

\msQuestion{11}
\begin{msscheme}{
  hline{3-5,7-13,15-16,18-19,21} = {2-Z}{0.4pt, msGrid},
  hline{6,14,17,20,22} = {0.4pt, msGrid},
  row{22} = {font=\bfseries},
}
Question & Scheme & Marks & Notes \\
11(a) & {\(\frac{\mathrm{d}^2x}{\mathrm{d}t^2} = -6\frac{\mathrm{d}x}{\mathrm{d}t} - 4\frac{\mathrm{d}y}{\mathrm{d}t}\)\\[4pt] \(\phantom{\frac{\mathrm{d}^2x}{\mathrm{d}t^2}} = -6\frac{\mathrm{d}x}{\mathrm{d}t} - 4(x - y)\)} & M1 & {Differentiates the first equation with respect to \(t\) and substitutes \(\frac{\mathrm{d}y}{\mathrm{d}t} = x - y\). Condone sign slips.} \\
& {\(4y = 40 - 6x - \frac{\mathrm{d}x}{\mathrm{d}t}\)\\[4pt] \(\Rightarrow \frac{\mathrm{d}^2x}{\mathrm{d}t^2} = -6\frac{\mathrm{d}x}{\mathrm{d}t} - 4x + 40 - 6x - \frac{\mathrm{d}x}{\mathrm{d}t}\)} & dM1 & {Eliminates \(y\) using the first equation, giving an equation in \(x\) and its derivatives only. Condone sign slips.} \\
& {\(\frac{\mathrm{d}^2x}{\mathrm{d}t^2} + 7\frac{\mathrm{d}x}{\mathrm{d}t} + 10x = 40\)} & A1~AG & {Correct substituted equation leading without error to the printed result; combined steps accepted.} \\
\nopagebreak & & [3] & \\
11(b) & {\(\lambda^2 + 7\lambda + 10 = 0 \Rightarrow \lambda = \ldots\)} & M1 & {Forms the auxiliary equation and solves it: factorising with product of constants 10, the formula, or completing the square. Calculator allowed.} \\
& {\(\lambda = -2, -5\)\\[4pt] \((x =)\ Ae^{-2t} + Be^{-5t}\)} & A1 & \\
& {Particular integral \(x = 4\)} & B1 & \\
& {\(x = Ae^{-2t} + Be^{-5t} + 4\)} & A1ft & {ft their complementary function of the form \(Ae^{\alpha t} + Be^{\beta t}\) + their constant particular integral.} \\
& {\(t = 0:\ \frac{\mathrm{d}x}{\mathrm{d}t} = 40 - 6(1) - 4(1) = 30\)\\[4pt] \(A + B + 4 = 1,\quad -2A - 5B = 30\)} & M1 & {Uses \(x\) at \(t = 0\) and \(\frac{\mathrm{d}x}{\mathrm{d}t}\) from the first equation in their general solution and its derivative, forming two equations in \(A\), \(B\). Condone arithmetic slips.} \\
& {\(\Rightarrow A = \ldots,\ B = \ldots\)} & dM1 & {Solves their two simultaneous equations for \(A\) and \(B\); calculator allowed.} \\
& {\(x = 5e^{-2t} - 8e^{-5t} + 4\)} & A1 & \\
\nopagebreak & & [7] & \\
11(c) & {\(\frac{\mathrm{d}x}{\mathrm{d}t} = -10e^{-2t} + 40e^{-5t} = 0\)\\[4pt] \(\Rightarrow e^{3t} = 4 \Rightarrow t = \ldots\)} & M1 & {Differentiates their \(x\), sets \(\frac{\mathrm{d}x}{\mathrm{d}t} = 0\) and uses logarithms to reach \(t = \ldots\). No verification of a maximum needed.} \\
& {\(t = \frac{1}{3}\ln 4\)} & A1 & {Or any exact equivalent, such as \(\frac{2}{3}\ln 2\).} \\
\nopagebreak & & [2] & \\
11(d) & {\(x = 3\times 4^{-\frac{2}{3}} + 4\)\\[4pt] \(40 - 6x - 4y = 0 \Rightarrow y = \ldots\)} & M1 & {Substitutes their \(t\) from (c) into their \(x\) and uses \(40 - 6x - 4y = 0\) to find \(y\).} \\
& {\(y = 4 - \frac{9}{2}\times 4^{-\frac{2}{3}}\ [\approx 2.214]\)\\[4pt] 221 ladybirds} & A1 & {Awrt 221 ladybirds; \(y \approx 2.21\) alone scores A0.} \\
\nopagebreak & & [2] & \\
11(e) & {For small \(x\), \(\frac{\mathrm{d}x}{\mathrm{d}t} \approx 40 - 4y\): eating rate \(4y\) does not depend on \(x\), so \(x\) can become negative} & E1 & {States that the \(-4y\) term does not depend on \(x\), so \(\frac{\mathrm{d}x}{\mathrm{d}t}\) can stay negative as \(x \to 0\): \(x\) negative, or aphids eaten when none remain.} \\
\nopagebreak & & [1] & \\
\nopagebreak \SetCell[c=2]{r} Total for Question 11 & & 15 & \\
\end{msscheme}

\msQuestion{12}
\begin{msscheme}{
  hline{3-7,9-10,12-14,16,18-19,21-24,26-28,30,32-33} = {2-Z}{0.4pt, msGrid},
  hline{8,11,15,17,20,25,29,31,34} = {0.4pt, msGrid},
  row{34} = {font=\bfseries},
}
Question & Scheme & Marks & Notes \\
12(a)(i) & {\(\frac{\mathrm{d}^2x}{\mathrm{d}t^2} = p\frac{\mathrm{d}x}{\mathrm{d}t} - a\frac{\mathrm{d}y}{\mathrm{d}t}\)\\[4pt] \(\phantom{\frac{\mathrm{d}^2x}{\mathrm{d}t^2}} = p\frac{\mathrm{d}x}{\mathrm{d}t} - a(qy - bx)\)} & M1 & {Differentiates the first equation with respect to \(t\) and substitutes the second for \(\frac{\mathrm{d}y}{\mathrm{d}t}\). Condone sign slips.} \\
& {\(ay = px - \frac{\mathrm{d}x}{\mathrm{d}t} \Rightarrow\)\\[4pt] \(\frac{\mathrm{d}^2x}{\mathrm{d}t^2} = p\frac{\mathrm{d}x}{\mathrm{d}t} - q\left(px - \frac{\mathrm{d}x}{\mathrm{d}t}\right) + abx\)} & dM1 & {Uses the first equation to replace \(y\), giving an equation in \(x\) and its derivatives only. Condone sign slips.} \\
& {\(\frac{\mathrm{d}^2x}{\mathrm{d}t^2} - (p+q)\frac{\mathrm{d}x}{\mathrm{d}t} + (pq-ab)x = 0\)} & A1 & {Correct substituted equation leading without error to the printed equation.} \\
& {\(x = e^{\lambda t} \Rightarrow \left(\lambda^2 - (p+q)\lambda + (pq-ab)\right)e^{\lambda t} = 0\)\\[4pt] \(e^{\lambda t} \neq 0 \Rightarrow \lambda^2 - (p+q)\lambda + (pq-ab) = 0\)} & M1 & {Substitutes \(x = e^{\lambda t}\) into the derived equation and cancels \(e^{\lambda t} \neq 0\). Stating the auxiliary equation alone scores M0.} \\
& {Distinct roots: \(e^{\lambda_1 t}\), \(e^{\lambda_2 t}\) both solutions\\[4pt] \(\Rightarrow x = Ae^{\lambda_1 t} + Be^{\lambda_2 t}\)} & A1~AG & {Concludes from two distinct roots that \(e^{\lambda_1 t}\) and \(e^{\lambda_2 t}\) are both solutions, combined with two arbitrary constants. The printed solution alone scores A0.} \\
\nopagebreak & & [5] & \\
12(a)(ii) & {\(\frac{\mathrm{d}x}{\mathrm{d}t} = \lambda_1Ae^{\lambda_1 t} + \lambda_2Be^{\lambda_2 t}\)\\[4pt] \(y = \frac1a\big[p(Ae^{\lambda_1 t} + Be^{\lambda_2 t})\)\\[4pt] \(\phantom{y = \frac1a\big[} - (\lambda_1Ae^{\lambda_1 t} + \lambda_2Be^{\lambda_2 t})\big]\)} & M1 & {Differentiates \(x\) and substitutes \(x\) and \(\frac{\mathrm{d}x}{\mathrm{d}t}\) into \(y = \frac1a\left(px - \frac{\mathrm{d}x}{\mathrm{d}t}\right)\) oe. Condone sign slips.} \\
& {\(y = \frac{p-\lambda_1}{a}Ae^{\lambda_1 t} + \frac{p-\lambda_2}{a}Be^{\lambda_2 t}\)} & A1~AG & {Correct substituted expression leading without error to the printed form.} \\
\nopagebreak & & [2] & \\
12(b)(i) & {\(\lambda^2 - 0.04\lambda - 0.0005 = 0\)\\[4pt] \(\Rightarrow \lambda = 0.05,\ -0.01\)} & B1 & \\
& {\(y = -\tfrac12Ae^{0.05t} + \tfrac32Be^{-0.01t}\)\\[4pt] \(A + B = x_0\), \(-\tfrac12A + \tfrac32B = y_0\)\\[4pt] \(\Rightarrow A = \ldots,\ B = \ldots\)} & M1 & {Uses \(t = 0\) in \(x\) and in their \(y\) of the form \(kAe^{0.05t} + mBe^{-0.01t}\) from (a)(ii), and solves the pair for \(A\) and \(B\).} \\
& {\(A = \tfrac14(3x_0 - 2y_0)\), \(B = \tfrac14(x_0 + 2y_0)\)\\[4pt] \(x = \tfrac14(3x_0 - 2y_0)e^{0.05t}\)\\[4pt] \(\phantom{x =} + \tfrac14(x_0 + 2y_0)e^{-0.01t}\)} & A1~AG & {Correct coefficients \(-\frac12\), \(\frac32\) in \(y\) and the pair solved without error to the printed result.} \\
\nopagebreak & & [3] & \\
12(b)(ii) & {\(y = -\tfrac12\cdot\tfrac14(3x_0 - 2y_0)e^{0.05t}\)\\[4pt] \(\phantom{y =} + \tfrac32\cdot\tfrac14(x_0 + 2y_0)e^{-0.01t}\)\\[4pt] \(\phantom{y} = -\tfrac18(3x_0 - 2y_0)e^{0.05t}\)\\[4pt] \(\phantom{y =} + \tfrac38(x_0 + 2y_0)e^{-0.01t}\)} & B1 & {Substitutes \(A\) and \(B\) from (b)(i) into \(y = -\frac12Ae^{0.05t} + \frac32Be^{-0.01t}\) and reaches the printed result without error.} \\
\nopagebreak & & [1] & \\
12(c)(i) & {\(x(25) = 50e^{1.25} + 350e^{-0.25}\)\\[4pt] \(y(25) = -25e^{1.25} + 525e^{-0.25}\)} & M1 & {Substitutes \(x_0 = 400\), \(y_0 = 500\) and \(t = 25\) into both results of (b).} \\
& {\(x \approx 447\), \(y \approx 322\)} & A1 & {awrt 447 and 322.} \\
\nopagebreak & & [2] & \\
12(c)(ii) & {\(50e^{0.05t} + 350e^{-0.01t} = -25e^{0.05t} + 525e^{-0.01t}\)} & M1 & {Equates their \(x\) and \(y\), each of the form \(Pe^{0.05t} + Qe^{-0.01t}\) with \(x_0\), \(y_0\) substituted.} \\
& {\(e^{0.06t} = \tfrac73\)} & A1 & {oe, such as \(75e^{0.06t} = 175\).} \\
& {\(0.06t = \ln\tfrac73 \Rightarrow t = \ldots\)} & dM1 & {Takes logarithms of their equation of the form \(e^{kt} = c\), \(c > 0\), to find \(t\).} \\
& {\(t \approx 14.1\) years} & A1 & {awrt 14.1.} \\
\nopagebreak & & [4] & \\
12(c)(iii) & {\(x = 50e^{0.05t} + 350e^{-0.01t} > 0\) for all \(t\)\\[4pt] so \(X\) never disappears} & E1 & {States both terms of \(x\) are positive for all \(t\) and concludes \(X\) never disappears.} \\
& {\(-25e^{0.05t} + 525e^{-0.01t} = 0\)\\[4pt] \(\Rightarrow e^{0.06t} = 21 \Rightarrow t = \ldots\)} & M1 & {Sets their \(y = 0\) and solves for \(t\).} \\
& {\(t \approx 50.7\): \(Y\) disappears after about 50.7 years} & A1 & {Concludes \(Y\) disappears, at awrt 50.7 years.} \\
\nopagebreak & & [3] & \\
12(d)(i) & {\(3x_0 = 2y_0\)} & B1 & {oe, such as \(y_0 = \frac32x_0\).} \\
\nopagebreak & & [1] & \\
12(d)(ii) & {\(x = \tfrac14(x_0 + 2y_0)e^{-0.01t} = x_0e^{-0.01t}\)\\[4pt] \(y = \tfrac38(x_0 + 2y_0)e^{-0.01t} = \tfrac32x_0e^{-0.01t}\)} & M1 & {Uses their relation from (d)(i) to remove the \(e^{0.05t}\) term from both \(x\) and \(y\).} \\
& {Both decrease exponentially to 0, with \(y : x = 3 : 2\)} & A1 & {Both populations decrease to zero, with constant ratio \(y : x = 3 : 2\) oe.} \\
\nopagebreak & & [2] & \\
\nopagebreak \SetCell[c=2]{r} Total for Question 12 & & 23 & \\
\end{msscheme}

\end{document}
